\scanpage{60}
Pour tt $i\in I$, $\unsure{i|\iota}(i)=i'$ est l'unique élément \uline{de $\Delta$} tel que $i'$ soit lié à \uline{tous les} él. de $I-\{i\}$. Donc la connaissance de $I$ implique déjà celle de $I'$. On appelle $I$ \uline{partie basique}. Si $I$ est telle, $I'$ est une partie basique, dite \uline{partie basique associée}. Base (ou partie basique ordonnée).

\deleted{\uline{\unsure{Prop}} La cat. des graphes cubiques}

Bibase : \deleted{bibase ordonnée \unsure{est un}} \inserted{él. de} $\mathfrak P_2(\mathfrak P_6(\Delta))$ formé de deux \deleted{\unsure{b}} bases associées.

\uline{Prop} a) La cat. \inserted{(pour les isos)} des graphes cubiques munis d'une base (resp. base ordonnée) est équivalente \inserted{can\textsuperscript{t}} à la cat. $\mathrm{Ens}_6$ des ens. à $6$ él (avec iso.) (resp. à la cat. ponctuelle)

b) La catégorie des gr. cub munis d'une bibase (resp. bibase ordonnée) est \inserted{can} équivalente à la catégorie \inserted{$(\mathrm{Ens}_6)\times(\mathrm{Ens}_2)$} des couples $(E,b)$ d'un ens. à six él. et d'un ens à $2$ él.

c) Soit $\Delta$ graphe cubique. L'ens des bases ordonnées de $\Delta$ est can. isom à
\[
\operatorname{Isom}_{\rm gr}(\overbrace{\Delta([1,6])}^{\Delta_0},\Delta),
\]
donc est un torseur à droite sous $W_0=\operatorname{Aut}\Delta_0\deleted{([1,6])}$, à gauche sous $W=\operatorname{Aut}(\Delta)$.

\uline{Th} Considérons $\Delta=\deleted{\Delta_0={}}\Delta(\unsure{\mathrm{II}|\mathbb I|I})\deleted{\unsure{\in\mathfrak P}}$, \deleted{\unsure{alors|Alors}} \inserted{les él. de $I$ dans $\Delta$ sont notés $\zeta_i$} avec sa base canonique \deleted{\unsure{ln}} $I$ et la base associée $I'$.

Les bases de $\Delta$ sont les parties \inserted{obtenues de la façon} suivantes.
% slanted
\margin{énoncer plutôt en termes de bibases \ldots}

(i) $I$ et $I'$ \deleted{\unsure{:}} (soit $2$ parties)

(ii) $\forall\ \deleted{\unsure{p\in\mathfrak P_2(\Delta)|J\in\mathfrak P_2(I)}}\ (i,j)\in I\times I-\operatorname{diag}$, $B^{ij}$ désigne la partie de $\Delta$ formée de $\{\zeta_j,\zeta'_j,(\zeta_{ik})_{k\in I-i-j}\}$ (soit $30$ parties)

(iii) $\forall K\in\mathfrak P_3(I)$, $B^{\deleted{\unsure{J}}K}$ désigne la partie \inserted{de $\Delta$} formée des $\{\zeta_i\ (i\in K),\ \deleted{\unsure{\cap}}\ \zeta_J\ (J\in\mathfrak P_2(I-K)\unsure{-|=}\mathfrak P_2(I))\}$ $(20)$

(iv) $\forall K\in\mathfrak P_3(I)$, $B'{}^K$ désigne la partie de $\Delta$
% archive sheet number: 59
